\sscp{Rodents and/or \it{Dasyuridae}}{C-03-05}
Another cognate set in languages of \ili{Aru} is glossed by \citet{Nivens2017}
as “mouse? Three-striped Dasyure?”.
If the second of these glosses is correct this would be another marsupial term.
For this cognate set \citet{Nivens2017} reconstructs Proto-\ili{Aru} **ŋiR(iu)wa.
Terms given by \citet{Nivens2017} as “ ‘tikus’ [‘mouse/rat’],
def[inition] not certain” are given as ``mouse/rat?''

\noindent{\wg\st{2pt}\begin{tabularx}{\linewidth}
{lllll>{\raggedright\arraybackslash}X}
\lsptoprule
%Language	&	ISO	&	Term	&	Gloss	&	Etymon	&	Source	\\	\midrule
%\endfirsthead
Language	&	ISO	&	Term	&	Gloss	&	Etymon	&	Source	\\	\midrule
%\endhead
\ili{Ujir}	&	ugj	&	diʤua,	&	\mc{2}{l}{Three-striped dasyure},		&		\cite{Nivens2017}\\	\hhline{~}
%	&		&		&	dasyure	&		&		\hp{\,}\\	
	&		&	daʤua	&	mouse/rat?	&		&	\hp{\,}	\\
%	&		&	daʤua	&	mouse/rat?	&		&	\hp{\,}	\\
\ili{Kompane}	&	kvp	&	dahia	&	mouse/rat?	&	**ŋiR(iu)wa	&	\cite{Nivens2017}	\\
West \ili{Kola}	&	kvv	&	daʤiwa/dayiwa	&	mouse/rat?	&	**ŋiR(iu)wa	&	\cite{Nivens2017}	\\
East \ili{Kola}	&	kvv	&	dahiwa/dehiwa	&	mouse/rat?	&	**ŋiR(iu)wa	&	\cite{Nivens2017}	\\
\ili{Lola}	&	lcd	&	dihó	&	mouse/rat?	&	**ŋiR(iu)wa	&	\cite{Nivens2017}	\\
\ili{Lorang}	&	lrn	&	ŋirugwa	&	mouse/rat?	&	**ŋiR(iu)wa	&	\cite{Nivens2017}	\\
\ili{Dobel}	&	kvo	&	ŋiro	&	mouse/rat?	&	**ŋiR(iu)wa	&	\cite{Nivens2017}	\\
\ili{Koba}	&	kpd	&	ŋarugwa	&	mouse/rat?	&	**ŋiR(iu)wa	&	\cite{Nivens2017}	\\
\ili{Barakai}	&	baj	&	deːróʔ	&	mouse/rat?	&	**ŋiR(iu)wa	&	\cite{Nivens2017}	\\
\ili{Manombai}	&	woo	&	derua	&	mouse/rat?	&	**ŋiR(iu)wa	&	\cite{Nivens2017}	\\
NE. \ili{Tarangan}	&	tre	&	ʤirua	&	mouse/rat	&	**ŋiR(iu)wa	&	\cite{Nivens2017}	\\
NW. \ili{Tarangan}	&	txn	&	ʤirua	&	mouse/rat	&	**ŋiR(iu)wa	&	\cite{Nivens2017}	\\
SW. \ili{Tarangan}	&	txn	&	nirua	&	mouse/rat	&	**ŋiR(iu)wa	&	\cite{Nivens2017}	\\
\lspbottomrule
\end{tabularx}\mbox{}\\}

\citet{Nivens2017} gives another cognate set which may designate
\it{Dasyuridae} and/or rodents, given below. For this set
\citet{Nivens2017} reconstructs Proto-South \ili{Aru} **kartá[βw](u).

\mbox{}\\\noindent{\wg\st{4pt}\begin{tabularx}{\linewidth}
{lllll>{\raggedright\arraybackslash}X}
\lsptoprule
%Language	&	ISO	&	Term	&	Gloss	&	Etymon	&	Source	\\	\midrule
%\endfirsthead
Language	&	ISO	&	Term	&	Gloss	&	Etymon	&	Source	\\	\midrule
%\endhead
\ili{Dobel}	&	kvo	&	ʔartáw	&	mouse/rat?	&	**kartá[βw](u)	&	\cite{Nivens2017}	\\
\ili{Barakai}	&	baj	&	kartá	&	mouse/rat?	&	**kartá[βw](u)	&	\cite{Nivens2017}	\\
\ili{Karey}	&	kyd	&	kartáo	&	mouse/rat?	&	**kartá[βw](u)	&	\cite{Nivens2017}	\\
\ili{Batuley}	&	bay	&	kartá	&	mouse/rat?	&	**kartá[βw](u)	&	\cite{Nivens2017}	\\
\ili{Mariri}	&	mqi	&	kartáw	&	mouse/rat?	&	**kartá[βw](u)	&	\cite{Nivens2017}	\\
NE. \ili{Tarangan}	&	tre	&	kartɔ	&	rat	&	**kartá[βw](u)	&	\cite{Nivens2017}	\\
SE. \ili{Tarangan}	&	tre	&	kartawtáw	&	mouse	&	**kartá[βw](u)	&	\cite{Nivens2017}	\\
NW. \ili{Tarangan}	&	txn	&	kartɔ	&	rat	&	**kartá[βw](u)	&	\cite{Nivens2017}	\\
SW. \ili{Tarangan}	&	txn	&	kartɔw	&	rat	&	**kartá[βw](u)	&	\cite{Nivens2017}	\\
\lspbottomrule
\end{tabularx}}